\chapter{Réels, intervalles et approximations}
\section*{Exercice 3.1 -- Classer}
Respectivement $\Z$, $\D$ car $7/8=0{,}875$, $\Q$ car $2/3$ n'a pas d'écriture décimale finie, $\N$ car $\sqrt{25}=5$, et $\R$ pour l'irrationnel $\sqrt3$. Pour ce dernier point, si $\sqrt3=p/q$ sous forme irréductible, alors $p^2=3q^2$. L'unicité de la décomposition en facteurs premiers impose que $3$ divise p ; écrire $p=3k$ donne $q^2=3k^2$, donc $3$ divise aussi q, contradiction.

\section*{Exercice 3.2 -- Intervalles}
$I\cap J=]0;2[$ ; $I\cup J=[-3;5]$ ; $\R\setminus I=]-\infty;-3[\cup[2;+\infty[$. La borne $2$ est exclue de $I$, donc incluse dans son complémentaire.

\section*{Exercice 3.3 -- Distance}
Le centre est $-2$ et le rayon $3$ : $x\in[-5;1]$. Pour le second intervalle, le centre est $7$ et le rayon $3$ : $|x-7|\leq3$.

\section*{Exercice 3.4 -- Encadrement}
$1{,}73^2=2{,}9929<3<3{,}0276=1{,}74^2$, donc $1{,}73<\sqrt3<1{,}74$. Pour l'arrondi, vérifier $1{,}735^2=3{,}010225>3$ : la racine est sous le milieu, donc l'arrondi est $1{,}73$. Le seul encadrement entre centièmes n'aurait pas suffi à choisir le plus proche.

\section*{Exercice 3.5 -- Irrationalité}
Si $1+\sqrt2$ était rationnel, lui soustraire le rationnel $1$ donnerait un rationnel. On obtiendrait que $\sqrt2$ est rationnel, contradiction. On utilise la stabilité de $\Q$ par soustraction.

\section*{Exercice 3.6 -- Chiffres significatifs}
On obtient $0{,}00457$ et $9{,}9\times10^4=99000$. Les pas d'arrondi sont respectivement $10^{-5}$ et $1000$ ; les erreurs absolues sont au plus $5\times10^{-6}$ et $500$. Ici les erreurs effectives valent $0{,}000003$ et $235$.
